\cotefile{157-2}
\scanpage{35}
\mspage{33}

le foncteur $(\sigma_M,\beta_M)$ précédent, i.e. qui donnent lieu à un \deleted{\unsure{diagramme}} \inserted{triangle} commutatif \unsure{(strict)}

\diagram{\begin{tikzcd}[ampersand replacement=\&]
 \& \underline{\mathrm{Ch}}_{\Sigma}(M) \arrow[d,"{(\sigma_M,\beta_M)}"] \\
 X \arrow[ur,"C"] \arrow[r,"{(F_0,F_1)}"] \& M\times M
\end{tikzcd}}

Mais si les foncteurs $F_0,F_1$ sont constants, de valeurs $a,b$ comme \unsure{tantôt|d'hab.}, i.e. $(F_0,F_1)$ constant de valeur $(a,b)$, alors \deleted{les} \inserted{la cat. des} foncteurs $C$ en question est \inserted{\deleted{\unsure{tau}} tautologiquement} isomorphe à celle des foncteurs de $\unsure{X}$ dans la fibre de $(\sigma_M,\beta_M)$ en $(a,b)$, d'où l'isomorphisme de catégories fondamental

\[
\boxed{\underline{\mathrm{Ch}}_{\Sigma_X}(X;a_X,b_X)\simeq
 \underline{\mathrm{Hom}}^{\unsure{!}}(X,\underline{\mathrm{Ch}}_{\Sigma}(M;a,b))}
\]

On va expliciter le contenu essentiel de cette formule en un

\uline{Scholie} \unsure{Soient} $M$ une catégorie, \unsure{avec|ou|soi} un ens $\Sigma$ de flèches de localisation, $a,b\in\mathrm{Ob}\,M$. Alors on peut trouver un \deleted{Ind} objet $\inserted{\underline{\underline K}=}K(a,b)$, \inserted{$K=(K_n)$} de $\mathrm{Ind}\,\mathrm{Cat}$\deleted{\unsure{, et}} (indexé par $\mathbb N$ mais peu importe au \unsure{fond|fait}) et un isom can, pour toute catégorie $X$ (isom fonctoriel en $X$)

\[
\begin{aligned}
\mathrm{Hom}_{\Sigma_X}(a_X,b_X)&\simeq
 \unsure{\Pi_0|\pi_0}\bigl(\underbrace{\underline{\mathrm{Hom}}(X,\underline{\underline K})}_{\substack{\Vert\,\unsure{\text{déf}}\\ \varinjlim\,\underline{\mathrm{Hom}}(X,K_n)}}\bigr)\\
 &\simeq\varinjlim\,\Pi_0\bigl(\underline{\mathrm{Hom}}(X,K_n)\bigr)\unsure{.}
\end{aligned}
\]

% archive sheet number: 34
\scanpage{36}
\mspage{34}

Je vais épiloguer sur ce résultat, sans me soucier pour l'instant de la façon particulière dont \deleted{$\underline{\underline K}$ a été \unsure{obt}} \unsure{les} \deleted{Ind-objets \inserted{\unsure{catégories} $K_n$ \unsure{(des Ind Cat)|(de Ind Cat)}} a été \unsure{obtenu}} \unsure{l'ind-|l'ind.}catégorie $K$ a été obtenue en termes de $M,\Sigma,a,b$.

\deleted{Il y a lieu} Comme les $\unsure{\Pi_0|\pi_0}\,\underline{\mathrm{Hom}}(X,K_n)=\mathrm{Hom}_{\mathrm{Cat\,hot}}(X,K_n)$ ne dépendent que de $X$, et des $K_n$ en tant qu'objets de Cat hot, \deleted{\unsure{je}} je vais interpréter $\underline{\underline K}$ comme un Ind-Objet de Cat hot, et écrire la formule fondamentale sous la forme

\[
\mathrm{Hom}_{\Sigma_X}(a_X,b_X)\simeq
 \mathrm{Hom}_{\deleted{\unsure{\mathrm{Ind\,Cat\,hot}}}\mathrm{Ind\text{-}Cat\text{-}hot}}(X,\underline{\underline K}).
\]

\unsure{Ainsi}, l'énoncé du Scholie implique deux choses :

a) Le \inserted{\unsure{contr.|contra.}} foncteur $X\unsure{\mapstochar\relbar|\longmapsto}\mathrm{Hom}_{\Sigma_X}(a_X,b_X)$ sur Cat. se factorise canoniquement par Cat hot, via un foncteur canonique

\diagram{\begin{tikzcd}[ampersand replacement=\&]
 \mathrm{Cat\,hot}^{\unsure{o|0}}\arrow[r,"{H=H_{a,b}^{M,\Sigma}}"] \& \mathrm{Ens}\mathpunct{;}
\end{tikzcd}}

b) Ce dernier \inserted{\unsure{contr.|contra.}} foncteur \deleted{\unsure{est}} sur Cat hot est \uline{ind-représentable}, par un ind-objet $\underline{\underline K}$ de Cat hot.

C'est cet ind-objet de Cat hot qui est défini à isom unique près par $M,\Sigma,a,b$, indépendamment de tout artifice de calcul pour expliciter les $\mathrm{Hom}_{\Sigma_X}(a_X,b_X)$.

Notons \deleted{\unsure{déjà}} aussi la conclusion suivante, nullement évidente à priori :

% archive sheet number: 35
\scanpage{37}
\mspage{35}

\deleted{\uline{Corollaire} Soit $f:X'\longrightarrow X^{\unsure{\prime|\;} }$ un foncteur qui soit une équivalence d'homotopie. \inserted{\unsure{Considérons} $M(X)\unsure{\relbar|\leftarrow}M(X')$, qui} Alors le foncteur correspondant $M(X^{\unsure{\prime|\;}})\Sigma_{X'}^{-1}\longrightarrow M$
\[
\mathrm{Hom}_{\Sigma_X}(a_X,b_X)\longrightarrow
\]}

Notons d'abord \deleted{que} (\unsure{ceci étant|c'était} implicite dans le Scholie) que toute flèche de Cat

\diagram{\begin{tikzcd}[ampersand replacement=\&]
 X \arrow[r,dash,"f"] \& Y
\end{tikzcd}}

définit

\diagram{\begin{tikzcd}[ampersand replacement=\&]
 M(Y)\arrow[r,"{M(f)}"] \& M(X)\quad\text{appliquant }\Sigma_Y\text{ dans }\Sigma_X,
\end{tikzcd}}

d'où

\[
 M(Y)\Sigma_Y^{-1}\overset{\mathrm{def}}{=}\mathbb D_{M,\Sigma}(Y)
 \longrightarrow M_X\Sigma_{\unsure{X|Y}}^{-1}=\mathbb D_{M,\Sigma}(X),
\]

\unsure{d'où}, pour tout $a,b\in\mathrm{Ob}\,M$, vu que $M(f)$ \deleted{applique} \inserted{transforme} $a_Y,\unsure{b_Y}$ en $a_X,b_X$ :

\[
 \boxed{\mathrm{Hom}_{\Sigma_Y}(a_Y,b_Y)\longrightarrow\mathrm{Hom}_{\Sigma_X}(a_X,b_X)}.
\]

Ceci posé, on a :

\uline{Corollaire} : Si $f$ est \unsure{une|un} \deleted{\unsure{équivalence d'}} \inserted{homotopisme}, alors quel que soient $a,b\in M$, la flèche précédente est un isomorphisme.

La question naturelle, c'est de savoir si la $=$ \deleted{\unsure{question}} conclusion est encore valable, en supposant seulement que $f$ soit une ``\uline{quasi-isomorphisme}'' (ou ``équivalence faible''), i.e que $f$ définisse un isom. dans la cat. \unsure{de|des} fractions Hot de Cat. (À ne pas confondre avec Cat hot,

% archive sheet number: 36
\scanpage{38}
\mspage{36}

nettement plus fine). S'il en était ainsi, le contr. foncteur

\[
 X\unsure{\relbar\joinrel\relbar|\longrightarrow}\mathrm{Hom}_{\Sigma_X}(a_X,b_X)
\]

se factoriserait non seulement par Cat hot, mais $=$ par la cat de fractions plus grossière Hot. Et dès lors se poserait la question si \deleted{\unsure{cet}} ce \unsure{contr.} foncteur hypothétique

\[
 \mathrm{Hot}^{\unsure{o|0}}\unsure{\relbar\joinrel\relbar|\longrightarrow}\mathrm{Ens}
\]

\unsure{est encore|est ou non} ind-représentable, \inserted{et} (de préférence\inserted{$^{*}$ \unsure{ind-repré.}} par l'ind-objet image de $\underline{\underline K}$ dans $\mathrm{Ind\,Hot}$ par

% slanted
\margin{\unsure{ce sera automatique} si on a ind-représentabilité comme foncteur sur Hot, cf ci-dessous.}

\[
 \mathrm{Ind}\,(\mathrm{Cat\,hot})\unsure{\relbar|\longrightarrow}\mathrm{Ind}\,(\mathrm{Hot}),
\]

déduit du \inserted{foncteur de \unsure{localis.}} $\mathrm{Cat\,hot}\unsure{\longrightarrow|\longleftarrow}\mathrm{Hot}$ \unsure{?)}. (Car cet ind-objet est un candidat tout trouvé.)

\deleted{On peut \deleted{\unsure{à}} \unsure{à toute} petite catégorie $K$, associer \inserted{fonctoriellement} une \uline{enveloppe injective} $\overline K$, et $K\unsure{\hookrightarrow|\relbar\joinrel\relbar}\overline K$ --- pour ce qui nous occupe il suffira que $\overline K$ soit une catégorie de \unsure{Thomason|Thomisson}-Kan. Ainsi un ind objet \inserted{$\underline{\underline K}=(K_n)$} de \unsure{Cat} (ici je \deleted{\unsure{préfère}} \unsure{utilise|utiliser} donc le fait que $\underline{\underline K}$ provient d'une vraie ind-catégorie) donne naissance à un ind-objet $\overline{\underline{\underline K}}=(\overline K_n)$, qui a $=$ image dans $\mathrm{Ind}\,(\mathrm{Hot})$, mais peut-être pas dans Ind Cat hot. On aura, pour un $X$ quelconque :}

% archive sheet number: 37
\scanpage{39}
\mspage{37}

Je vais \unsure{faire} quelques rappels catégoriques, généralisant la situation du foncteur canonique de localisation

\diagram{\begin{tikzcd}[ampersand replacement=\&]
 \mathrm{Cat\,hot}\arrow[r,"\varphi"'] \& \deleted{\unsure{C}}\mathrm{Hot},
\end{tikzcd}}

dont on sait qu'il admet un adjoint à droite

\diagram{\begin{tikzcd}[ampersand replacement=\&]
 \mathrm{Cat\,hot} \& \mathrm{Hot}\arrow[l,"\psi"'] \&[-1em] ,\quad
 \mathrm{Hom}_{\mathrm{Hot}}(\varphi X,Y)\simeq\mathrm{Hom}_{\mathrm{Cat\,hot}}(X,\psi Y)
\end{tikzcd}}

où $\psi$ est le foncteur ``enveloppe de Kan'' ou ``représentant de Kan''. On va donc considérer la situation plus générale

\diagram{\begin{tikzcd}[ampersand replacement=\&]
 C\arrow[r,shift right,"\varphi"'] \& C'\arrow[l,shift right,"\psi"'] \&[-1em]
 \mathrm{Hom}_{C'}(\varphi X,X')\simeq\mathrm{Hom}_C(X,\psi X')
\end{tikzcd}}

i.e. $\varphi,\psi$ couple de foncteurs adjoints.

On suppose de plus

\[
 \left\{\begin{aligned}
 &\varphi\ \text{un foncteur localisation, i.e. }\psi\ \text{pleinement fidèle},\\
 &\text{i.e.}\qquad\deleted{\unsure{\mathrm{id}_{C'}\longrightarrow}}
  \varphi\psi\xrightarrow{\sim}\mathrm{id}_{C'} .
 \end{aligned}\right.
\]

Soit $W=W_{\varphi}\subset\mathrm{Fl}\,C$, \unsure{syst} défini par $\varphi$: $W_\varphi=\{u\in\mathrm{Fl}\,C\mid\varphi(u)\ \mathrm{iso}\}$. Voici la proposition que j'avais oubliée, et que j'ai eu du mal à réconstituer, hélas !

\uline{Proposition 1} Soit $K\in\mathrm{Ob}\,C$. Les conditions suivantes \inserted{sur $K$} sont équivalentes :

(i) Le foncteur

\[
 X\longrightarrow K(X)\overset{\mathrm{def}}{=}\mathrm{Hom}_C(X,K)
 \qquad C^{\unsure{o|0}}\unsure{\relbar|\longrightarrow}\mathrm{Ens}
\]

transforme flèches de $W$ (``quasi-isomorphismes) en bijections

(i$'$) \unsure{Même|Mêmes} \unsure{condition|conditions} pour les flèches d'adjonction

\[
 X\longrightarrow\psi\varphi X\qquad(\forall X\in\mathrm{Ob}\,C)
\]

(i$\unsure{''}$) Même condition pour la flèche d'adjonction

\[
 K\longrightarrow\psi\varphi K
\]

\margin{N.B. Il suffit $=$ que l'élém $\mathrm{id}_K$ de $F_K(K)$ provienne de $\unsure{F_K}(\psi\varphi K)$ i.e. (i$'''$) Il existe une rétraction de $\psi\varphi(X)$ sur $X$ (\uline{NB} Celle-ci sera \unsure{même|bien} un iso)}

(ii) $K$ appartient à l'image essentielle de $\psi$

(ii$'$) $K\longrightarrow\psi\varphi K$ (adjonction) est un iso.

(ii$''$) L'hom fonctoriel en $X$

\[
 \mathrm{Hom}_C(X,K)\longrightarrow\mathrm{Hom}_{C'}(\varphi X,\varphi K)
\]

est un iso.

\uline{Dém} On a d'abord comme quasi-tautologies

\diagram{\begin{tikzcd}[ampersand replacement=\&,column sep=small,row sep=small]
 (\mathrm{ii}\unsure{'})\arrow[r,Rightarrow] \& (\mathrm{ii}\unsure{\vphantom{' }|'})\arrow[r,Rightarrow] \&(\mathrm{ii}'')\arrow[d,Rightarrow]\\
 \& \&(\mathrm{ii}')
\end{tikzcd}}

i.e. l'équivalence des trois dernières conditions. En effet (ii$'$) $\Longrightarrow$ (ii) tautologique, (ii) $\Longrightarrow$ (ii$''$) car si $K=\psi(K')$ on a

% archive sheet number: 38
\scanpage{40}
\mspage{38}

par adjonction

\[
 \mathrm{Hom}_C(X,\underset{\substack{\Vert\\\psi(K')}}{K})
 =\mathrm{Hom}_{\unsure{C|C'}}(X,\psi(K'))
 \overset{\mathrm{adj.}}{\simeq}\mathrm{Hom}_{C'}(\varphi X,\deleted{\unsure{\varphi\psi}}(K'))
\]

\deleted{or $\varphi\psi(K')\xrightarrow{\sim}K'$ (adj, et le fait que $\psi$ soit pl. fid), donc} \inserted{or $\varphi(K)=\varphi\psi(K')\simeq K'$ (adj, et $\psi$ pl. fid)}

\diagram{\begin{tikzcd}[ampersand replacement=\&,column sep=small]
 \mathrm{Hom}_C(X,K)\arrow[rr,"\sim"]\arrow[dr,"{\varphi_{X,K}}"] \& \& \mathrm{Hom}_{C}(\varphi X,K')\\
 \& \mathrm{Hom}_{C'}(\varphi X,\underbrace{\varphi K}_{\substack{\deleted{\unsure{\psi\varphi K'}}\\ \Vert\,\varphi\psi(K')}})\arrow[ur,"\sim","\mathrm{adj}"'] \&
\end{tikzcd}}

\inserted{qu'on complète en}

ce qui prouve que $\varphi_{X,K}$ est iso.

(ii$''$) $\longrightarrow$ (ii$'$) \unsure{Car|car} \deleted{\unsure{la}} flèche

\diagram{\begin{tikzcd}[ampersand replacement=\&]
 \mathrm{Hom}_C(X,K)\arrow[r,"\alpha_X"]\arrow[dr,dashed] \&
 \mathrm{Hom}_{C'}(\varphi X,\varphi K)\arrow[d,dash,"\wr\ \mathrm{adj}"]\\
 \& \mathrm{Hom}_C(X,\psi\varphi K)
\end{tikzcd}}

est déduite de la flèche d'adjonction

\[
 K\longrightarrow\psi\varphi K
\]

en \unsure{prenant|posant} $\mathrm{Hom}_C(X,-)$, donc $\alpha_X$ un iso pour tout $X$ ssi cette flèche d'adjonction est iso. Donc

\[
 \mathrm{ii}\Longleftrightarrow\mathrm{ii}'\Longleftrightarrow\mathrm{ii}''
\]

D'autre part,

\[
 \mathrm{ii}''\Longrightarrow\mathrm{i}\Longrightarrow\mathrm{i}'\Longrightarrow\mathrm{i}''
\]

\margin{Car la flèche d'adjonction \deleted{$\unsure{X\to\psi\varphi X}$} $X\to\psi\varphi X$ \unsure{dans $C$ est|donc est} transformée par $\varphi$ en $\varphi X\to\underbrace{\varphi\psi}\varphi(X)$ dont le composé avec $(\varphi\psi)\varphi X\to\varphi X$}

sont tautologiques, \deleted{il reste à prouver $\mathrm{i}''\Longrightarrow\mathrm{ii}'$,} \unsure{en tenant compte|en tant} pour (i) $\Longrightarrow$ \unsure{(i)|(i$'$)} que \deleted{i.e. en supposant que pour la flèche} les flèches d'adjonction $X\xrightarrow{u}\psi\varphi X$ de $C$ sont bien dans $W$, car on a un triangle commutatif

\diagram{\begin{tikzcd}[ampersand replacement=\&]
 \varphi(X)\arrow[r,"{\varphi(u)}"]\arrow[dr,swap,"{\mathrm{id}_{\varphi(X)}}"] \&
 \varphi(\psi\varphi X)=(\varphi\psi)(\varphi X)\arrow[d,"{\wr\ \mathrm{adj.}\ \varphi\psi\to\mathrm{id}\quad\text{qui est iso car }\psi\ \text{pl. fid}}"]\\
 \& \varphi(X)
\end{tikzcd}}

donc $\varphi(u)$ bien iso. Reste à prouver $\mathrm{i}''\Longrightarrow\mathrm{ii}'$, i.e. on suppose

% archive sheet number: 39
\scanpage{41}
\mspage{39}

que le foncteur $X\unsure{\mapstochar\relbar|\longmapsto}\mathrm{Hom}(X,K)$ transforme le morphisme d'adjonction particulier \unsure{$K$}

\[
 K\xrightarrow{i}\psi\varphi K,
\]

en

\[
 \mathrm{Hom}\bigl(\psi\varphi(K),K\bigr)
 \xrightarrow{u\unsure{\relbar|\mapsto}u\circ i}\mathrm{Hom}(K,K)
\]

qui est une bijection. \deleted{\unsure{Donc}} $\mathrm{id}_K$ est dans l'image, soit $p$ un élément du premier ens qui lui donne naissance

\diagram{\begin{tikzcd}[ampersand replacement=\&]
 K\arrow[r,hook,"i"] \& \psi\underbrace{\varphi(K)}_{\Vert\,K'}\arrow[l,bend left=22,"p"] \&[-1em] \text{donc }p\circ i=\mathrm{id}_K
\end{tikzcd}}

\unsure{Considérons|considérons}

\[
 ip=q:\psi(K')\longrightarrow\psi(K')\qquad(\mathrm{NB}\ q^2=q)
\]

comme $\psi$ est pl. fid. il existe un unique

\[
 q_0:K'\longrightarrow K'\qquad(K'\overset{\mathrm{def}}{=}\varphi(K))
\]

tel que l'on ait

\[
 q=\psi(q_0).
\]

La relation $pi=\mathrm{id}_K$ donne, en appliquant $\varphi$

\[
 \varphi(p)\varphi(i)\inserted{\unsure{\mathrm{fid.}}}=\mathrm{id},\quad\unsure{\text{or}}
\]

donc comme $\varphi(i)$ est iso, (c'est \unsure{évident comm. du|en faisant comm. le} triangle ci-contre),

\diagram{\begin{tikzcd}[ampersand replacement=\&]
 \varphi(K)\arrow[r,"\sim"]\arrow[dr,dash,swap,"{\sim\,\mathrm{id}}"] \& \varphi(\psi\varphi(K))=(\varphi\psi)(\varphi K)\arrow[d,"{\wr\ \mathrm{adj.}}"]\\
 \&\varphi K
\end{tikzcd}}

$\varphi(p)$ l'est. Donc $\varphi(q)=\varphi(i)\varphi(p)$ l'est, puisque $\varphi(i)$ et $\varphi(i)$ le sont. Mais

\[
 \varphi(q)=\varphi\psi(q_0)\xrightarrow[\ ]{\substack{\mathrm{adj}\ \varphi\psi\simeq\mathrm{id}\\ \sim}}q_0,
\]

donc $q_0$ est un iso, donc aussi $q=\psi q_0$. Comme $q^2=q$, cela implique $q=\mathrm{id}_{\unsure{K|{?}}}$. Donc $i$ et $p$ sont des iso. inv. l'un de l'autre, ce qui établit \unsure{(ii)$''$|(ii$'$)} et achève la dém.

\uline{Corollaire} : Considérons deux foncteurs $C^{\unsure{o|0}}\xrightarrow{H}\mathrm{Ens}$, et $C'{}^o\xrightarrow{H'}\mathrm{Ens}$, tels que $H^{\deleted{\unsure{\prime}}}\simeq H'\circ\varphi$, \deleted{\unsure{alors}}. Pour que $H$ soit représentable, il faut et \unsure{s} que $H'$ le soit, et dans ce cas, l'objet $K$ de $C$

% archive sheet number: 40
\scanpage{42}
\mspage{40}

qui représente $H$ satisfait aux conditions équivalentes de la prop. Si $K,K'$ les représentent, on a des iso. can. $K'\unsure{=|\simeq}\varphi K,\ K\unsure{=|\simeq}\psi K'$.

Si $H'$ est représentable par $K'\in\mathrm{ob}\,C'$, on aura donc

\[
 H'\circ\varphi\,(X)\simeq\mathrm{Hom}_{C'}(\varphi X,K')\overset{\mathrm{adj}}{\simeq}
 \mathrm{Hom}_C(X,\psi(K'))
\]

donc $H$ est représentable par $K=\psi(K')$, satisfaisant la condition (ii) de la prop. Inversement, si \deleted{\unsure{l'on a}} $H$ est représentable par un objet $K$, donc $K$ satisfait à la Condition (i), donc \deleted{il est \unsure{isom.} à un $\psi(K')$}, \inserted{par (ii$''$) on a} \deleted{d'où}

\[
 \deleted{\mathrm{Hom}_C(X,K)\simeq\mathrm{Hom}_{\unsure{C'|C}}(X,\psi(K'))\overset{\mathrm{adj}}{\simeq}}.
\]

\[
 \mathrm{Hom}_C(X,K)\simeq\mathrm{Hom}_{C'}(\varphi X,\varphi K),
\]

donc le foncteur $H':C'{}^{\unsure{o|0}}\unsure{\relbar|\longrightarrow}\mathrm{Ens}$ qui factorise $H$ s'explicite comme $X'\to\mathrm{Hom}_{C'}(X',\varphi K)$, \deleted{\unsure{et}} donc est représentable par $\varphi K$.

Passons \deleted{\unsure{aux}} aux ind-objets. \deleted{Soit} Les foncteurs $\varphi,\psi$

\deleted{$(K_i)_{i\in I}$ un ind-objet de $C$, et soit $H$ le foncteur qui le représente. Supposons}

passent aux Ind-objets, d'où des foncteurs

\diagram{\begin{tikzcd}[ampersand replacement=\&]
 \mathrm{Ind}\,C\arrow[r,shift left,{|->},"{\Psi=\mathrm{Ind}\,\psi}"]
 \arrow[r,shift right,"{\Phi=\mathrm{Ind}\,\varphi}"'] \& \mathrm{Ind}\,C'
\end{tikzcd}}

On aura encore la formule d'adjonction pour $X=(X_i)_{i\in I}$ de $\mathrm{Ind}\,C$, $X'=(X'_{i'})_{i'\in I'}$ de $\mathrm{Ind}\,C'$, car

\[
\begin{aligned}
 \mathrm{Hom}(X,\overset{\deleted{\unsure{\mathrm{Hom}}}}{\Psi}X')&
 \simeq\varprojlim_i\varinjlim_{i'}\mathrm{Hom}(X_i,\psi(X'_j))\\
 &\simeq\varprojlim_i\varinjlim_{i'}\mathrm{Hom}(\varphi X_i,X'_j)
 \simeq\mathrm{Hom}(\deleted{\unsure{\varphi}}\Phi X,\deleted{\unsure{\psi}}X').
\end{aligned}
\]

% archive sheet number: 41
\scanpage{43}
\mspage{41}

On prouve de même que $\Psi$ est pl. fidèle, car si $X'=(X'_i)_{i\in I}$, $Y'=(Y'_j)_{j\in J}$ sont deux objets de $\mathrm{Ind}\,C'$, on aura

\diagram{\begin{tikzcd}[ampersand replacement=\&,column sep=small]
 \mathrm{Hom}(X',Y')\arrow[d,"{\Psi_{X',Y'}}"] \& {}={} \& \varprojlim_i\varinjlim_j \& \mathrm{Hom}_{C'}(X'_i,Y'_j)\arrow[d,"{\psi_{X'_i,Y'_j}}"]\\
 \mathrm{Hom}(\Psi X',\Psi Y') \& {}={} \& \varprojlim_i\varinjlim_j \& \mathrm{Hom}_{C}(\psi(X'_{\unsure{i|j}}),\psi(Y'_j)),
\end{tikzcd}}

\deleted{\unsure{la première}} donc $\Psi_{X',Y'}$ est iso car les $\psi_{X'_i,Y'_j}$ le sont.

Donc la proposition 1 peut s'appliquer à la situation \deleted{$\unsure{\Phi,\Psi}$} $\Phi,\Psi$ entre Ind-objets. Je vais l'expliciter sous la forme la plus commode :

\uline{Proposition 2} Soit $K=(K_i)_{i\in I}$ un ind-objet de $C$, \deleted{\unsure{Conditions}, \unsure{équivalentes}} \inserted{et soit $H$ le foncteur}

\[
 \left\{\begin{aligned}
 H:\ X&\unsure{\relbar|\longrightarrow}\underbrace{\mathrm{Hom}_C(X,K)}_{K(X)}
 \overset{\mathrm{def}}{=}\varinjlim_i\mathrm{Hom}_C(X,K_i)\\
 C^{\unsure{o|0}}&\unsure{\relbar|\longrightarrow}\mathrm{Ens}
 \end{aligned}\right.
\]

qu'il ind-représente. Conditions équivalentes

(i) Le foncteur $H$ se factorise (à isom près) par $C'$ via $\varphi$, i.e transforme les $u\in W_\varphi$ en iso.

(i$\unsure{'}$) Même condition pour les flèches d'adjonction $X\to\psi\varphi(X)$ (lesquelles sont dans $W$).

\inserted{(i$''$) Même condition pour la flèche d'adjonction $K\to\Psi\Phi K$.}

(ii) $K$ appartient à l'image essentielle de $\mathrm{Ind}\,C'$ par $\Psi$.

(ii$'$) $K\longrightarrow\Psi\Phi(K)$ est un \uline{iso} de $\mathrm{Ind}\,C$

(ii$''$) L'hom fonctoriel

\[
 K(X)=\mathrm{Hom}_C(X,K)\longrightarrow\mathrm{Hom}_{C'}(\varphi(X),\varphi(K))
\]

est un iso.

\uline{Dém} L'\unsure{équivalence|équi.} de ii, ii$'$, ii$''$ est la prop. 1 appliquée à $\Phi$ et $\Psi$, \deleted{Et ii$''$ implique} à cela près qu'il faut \unsure{noter|montrer} que ii$''$ reste \inserted{évidemment} valable si on y remplace $X$ par un ind-objet $(X_j)_{j\in J}$ quelconque. On

% archive sheet number: 42
\scanpage{44}
\mspage{42.}

\[
 \mathrm{ii}''\Longrightarrow\mathrm{i}\Longrightarrow\mathrm{i}'\deleted{\unsure{\relbar}}
\]

tautologiquement, il reste à prouver que i$'$ implique \inserted{i$''$ et i$''$} (ii$'$). On ne peut appliquer directement la prop. 1, car l'hyp. \deleted{\unsure{n'est}} \inserted{(i$'$)} ne porte que sur les objets de $\mathrm{Ind}\,C$ qui proviennent de $C$, et il n'est pas évident à priori qu'elle passe aux objets quelconques.

Mais cela l'est dans le cas d'un hom de ind-objets sur \unsure{les} $(X_i),(Y_i)$ sur \unsure{le|la} $=$ \unsure{ens|cat} d'indices, \deleted{\unsure{qui est un}} \inserted{donné par} \deleted{\unsure{hom} $X_i\unsure{\relbar|\to}\unsure{Z|Y_i}$} un \deleted{\unsure{syst.}} hom ordinaire (\unsure{terme à terme}) de \deleted{\unsure{inductifs}} systèmes inductifs sur $I$, $X_i\to\unsure{Z_i|Y_i}$ : \deleted{\unsure{Si les}} si les $X_i\unsure{\relbar|\to}Y_i$ sont \unsure{dans|des} $W$, l'hyp. (i) implique que $H(Y)\to H(X)$ est iso, et de $=$ pour i$'$ pour

\[
 \underset{\substack{\Vert\\(X_i)}}{X}\longrightarrow
 \underset{\substack{\Vert\\(\psi\varphi(X_i))_i}}{\Psi\Phi(X)} ,
\]

où l'hom. d'adjonction est donné terme à terme par les hom d'adjonction $X_i\longrightarrow\psi\varphi(X_i)$. Donc l'hyp. i$'$ \deleted{de prop. \unsure{2}} ici implique l'hyp. i$'$ de prop 1 pour $\Phi$ et $\Psi$, donc aussi \unsure{(ii$'$) ici} (par prop. 1), et on gagne.

\figure{\begin{tikzpicture}[x=1cm,y=1cm]\draw (0,0.28)--(0.15,0);\end{tikzpicture}}

Revenons à la situation où $C=\mathrm{Cat\,hot}$, $C'=\mathrm{Hot}$, $\varphi$ le foncteur localisation. Si $K\in\mathrm{Ob}\,\mathrm{Cat\,hot}$, alors $\psi\varphi(K)$ s'appelle \uline{l'enveloppe de \inserted{Thomason-}\unsure{Kan}} de $K$ --- si $K$ est une catégorie, son enveloppe $\overline K$ \deleted{et le} et la \unsure{flèche|flèches} $K\unsure{\relbar|\to}\overline K$ ne sont définis que dans la catégorie Cat hot, i.e. $\overline K$ n'est défini qu'à ``homotopie près'',

% archive sheet number: 43
\scanpage{45}
\mspage{43}

Mais \unsure{on} peut toujours \uline{choisir} \deleted{$\unsure{K,K}$} un \deleted{$\unsure{\overline K}$} objet \inserted{$\overline K$} de Cat, et une flèche

\[
 i:K\longrightarrow\overline K\qquad\text{\uline{dans Cat}}
\]

pas seulement \deleted{\unsure{dans}} Cat hot, de telle façon de plus que $i$ soit une ``\uline{cofibration} de Thomason''. (et en particulier une \unsure{imm ouverte|imm. ouverte} ?). Il \deleted{s'ensuit} en résulte que si on a un système inductif

\[
 K=(K_n)_{n\in\mathbb N}
\]

indexé par $\mathbb N$, et si \deleted{\unsure{on}} on choisit pour chaque $n$ une enveloppe de Thomason-Kan

\[
 K_n\overset{\alpha_n}{\hookrightarrow}\overline K_n
\]

de cette façon, on peut trouver des \unsure{morphismes|flèches} de transition

\[
 \overline K_n\overset{\overline{\unsure{i_n|j_n}}}{\unsure{\relbar|\longrightarrow}}\overline K_{n+1}
\]

\unsure{tels} que les carrés

\diagram{\begin{tikzcd}[ampersand replacement=\&]
 K_n\arrow[r,hook,"{\alpha_n}"]\arrow[d,swap,"{i_n}"] \& \overline K_n\arrow[d,"{\overline{i_n}}"]\\
 K_{n+1}\arrow[r,hook,"{\alpha_{n+1}}"] \& \overline K_{n+1},
\end{tikzcd}}

soient commutatifs ; de telle sorte que l'on ait donc un hom de \uline{systèmes inductifs de Cat}

\[
 K\longrightarrow\overline K
\]

et a fortiori un hom de Ind-objet de Cat. Ceci posé, les conditions \inserted{(i),} (ii$'$) de la prop 2 montrent que le foncteur ind-représenté par $K$ \uline{dans $C=\mathrm{Cat\,hot}$} \deleted{\unsure{est}} transforme quasi-isomorphismes

% archive sheet number: 44
\scanpage{46}
\mspage{44}

ou \unsure{du moins|disons} (équivalences faibles) en bijections, \uline{ssi} l'hom précédent est un isom. de ind-objets de Ind Cat hot.

Ce sera donc cela la condition naturelle à imposer aux \deleted{\unsure{catég}} ind-catégories de chemins

\[
 \underline{\underline{\mathrm{Ch}}}_{\Sigma}(M,\unsure{a,b|a;b})
 \qquad(a,b\in\mathrm{Ob}\,M)
\]

provenant d'une ``catégorie de modèles'' $M,\Sigma$.

Bien sûr, disposant d'un système inductif explicite de catégories (pas seulement d'objets de Cat hot)

\[
 K_n=\underline{\mathrm{Ch}}_{\Sigma,n}(M;a,b),
\]

il est tentant de passer à la limite inductive

\[
 K_\infty=\varinjlim_n\underline{\mathrm{Ch}}_{\Sigma,\deleted{\unsure{n|\infty}}}(M;a,b)
 \overset{\mathrm{def}}{=}\underline{\mathrm{Ch}}_\Sigma(M;a,b),
\]

et ceci d'autant plus que les foncteurs de transition sont des monomorphismes et pleinement fidèles, de sorte que \deleted{$K_{\unsure{\infty|n}}$ est} les $K_n$ s'identifient à une suite de souscatégories pleines de $K_\infty$, dont la réunion est $K_\infty$. Mais d'autre part je ne sais donner aucune interprétation \unsure{convaincante} de $K_\infty$, \unsure{comme|connexe} \deleted{de} de celle du ind-objet \inserted{de Cat hot} défini par $\underline{\underline K}=(K_n)_{n\in\mathbb N}$, qui ind-représente le foncteur remarquable

\[
 H=H_{\unsure{M,\Sigma|M},a,b}:\mathrm{Cat\,hot}^{\unsure{o|0}}
 \unsure{\relbar\joinrel\relbar|\longrightarrow}\mathrm{Ens}
\]

factorisant le foncteur

% archive sheet number: 45
\scanpage{47}
\mspage{45}

\[
 X\longmapsto H_{M;x,y}(X)\overset{\mathrm{def}}{=}
 \mathrm{Hom}_{M(X)\Sigma_X^{-1}}(\deleted{\unsure{a}}\unsure{x|a}_X,\deleted{\unsure{b}}\unsure{y|b}_X),
\]

\deleted{\unsure{Non seulement}} \inserted{par la formule donc}

\[
 H(X)=\varinjlim_n
 \underbrace{\unsure{\Pi_0|\pi_0}\,\underline{\mathrm{Hom}}(X,K_n)}_{\substack{\Vert\\\mathrm{Hom}_{\mathrm{Cat\,hot}}(X,K_n)}} .
\]

Je ne sais donner d'interprétation de ce type ni de la catégorie $K_\infty$ elle-même, ni de son image dans Cat hot, ni de son image dans Hot.

Il est vrai que si je \unsure{restreint|restreins} le foncteur $H$ à \deleted{\unsure{ceux}} la sous-catégorie pleine de (Cat) formée des catégories \uline{finies} ($\mathrm{Ob}\,X$ et $\mathrm{Fl}\,X$ finis), alors on aura

\[
 \underline{\mathrm{Hom}}^{!}(X,K_\infty)\unsure{=|\simeq}\underline{\mathrm{Hom}}(X,K_\infty),
\]

et l'isom. précédent s'interprète comme

\[
 \underset{\substack{\Vert\,\mathrm{def}\\
 \mathrm{Hom}^{\unsure{\prime|\;}}_{M(X)\Sigma_X^{-1}}(\unsure{a|x}_X,\unsure{b|y}_Y)}}{H(X)}
 \simeq\Pi_0\,\underline{\mathrm{Hom}}(X,K_\infty)\overset{\text{déf.}}{\unsure{=|\simeq}}
 \mathrm{Hom}_{\mathrm{Cat\,hot}}(X,K_\infty).
\]

\deleted{\unsure{Ce}} Cela prouve du moins que la connaissance de $K_\infty$ est utile. Mais cette formule, pour $X$ restreint aux catégories \uline{finies}, ne détermine pas $K_\infty$ \deleted{``à homotopie} comme objet de Cat hot, à isom \inserted{unique} près dans cette catégorie --- \uuline{sauf} dans le cas \deleted{\unsure{bien}} \deleted{\unsure{douteux}} où $K_\infty$ est elle-$=$ \deleted{\unsure{don}} homotope à une \unsure{catégo}

% archive sheet number: 46
\scanpage{48}
\mspage{46}

rie finie. (Je ne me rends pas bien compte si \unsure{sa|ça} arrive très souvent\ldots\ \deleted{\unsure{)}} Mais on ne peut fonder une théorie générale des dérivateurs et des catégories de modèles ``dérivables'', sur une hypothèse de ce genre.)

Mais supposons que le foncteur $H$ envisagé se factorise même par \unsure{Hot}, i.e. qu'il transforme quasi-iso. en bijections. On a vu que cela équivaut à la condition qu'il puisse se représenter par un système inductif sur $\mathbb N$ formé de catégories de Thomason-Kan.

Il doit être facile de plus, si on y tient, de s'arranger pour que les foncteurs de transition soient des mono et pl. fid., pour avoir une $\underline{\lim}$ \unsure{sympa}. Ce sera bien \unsure{sur|sûr} une catégorie de Thomason-Kan.

Cette catégorie est-elle déterminée par le foncteur $H:\mathrm{Cat\,hot}^{\unsure{o|0}}\unsure{\relbar|\to}\mathrm{Ens}$, à iso unique près \unsure{!|à} dans Cat hot ? Bien sûr, il n'est pas question qu'elle représente le foncteur $H$, \uline{sauf} si \inserted{celui-ci est bel et bien représentable, i.e.} le \deleted{système inductif} \inserted{ind objet} $\underline{\underline K}$ de Cat hot est essentiellement constant (ce qui ne doit pratiquement jamais être le cas). Notons déjà

% archive sheet number: 47
\scanpage{49}
\mspage{47}

que si on \unsure{un|a un} hom. de ind-objets \uline{de Cat}, \unsure{ou.|ou,} ce qui revient maintenant au même, un hom \deleted{de} Cat

\[
K_\infty \xrightarrow{u_\infty} K'_\infty
\]

tel que $\forall n,\quad u_\infty(K_n)\subset K'_{\deleted{\unsure{\tau(n)}}n'}$ pour $n'=\tau(n)$ convenable, \deleted{et si \unsure{cet hom.|cet} \unsure{est tel|est} que} et si on suppose que l'hom de ind-objets induit un isom. pour les foncteurs ind-représentés \uline{sur Cat hot} (p. ex. si $\underline{\underline K}$ et $\underline{\underline {K'}}$ \inserted{ind-}représentent le $=$ foncteur $H$, et l'hom de ind-objets est celui qui correspond à $\mathrm{id}_H$), alors $u_\infty$ est forcément un isom. \uline{dans Hot}, donc \unsure{aussi} \inserted{(sauf erreur)} dans Cat hot puisque $K_\infty,\ K'_\infty$ sont des catégories de Thomason-Kan. Mais ce qui n'est pas évident, c'est l'existence d'un hom $u_\infty$ comme dessus, et s'il existe, qu'il soit unique \inserted{(ou du moins canonique ?)} comme flèche de Cat hot.

Pour l'existence d'un $\unsure{\dot u_\infty|u_\infty}$, on est ramené à prouver que si on a un diagramme à flèches pleines

\diagram{\begin{tikzcd}[ampersand replacement=\&,column sep=large]
K_n \arrow[r,"i_n","{\mathrm{cofib.}}"'] \arrow[d,"u_n"] \& K_{n+1} \arrow[dr,dashed,"u_{n+1}"] \& {}\\
K'_{\tau(n)} \arrow[rr,"{\text{transition}}"' ,"{\unsure{\scriptscriptstyle\blacktriangleright|\cdot}}" marking] \& {} \& K'_N
\end{tikzcd}},

et si une flèche $u_{n+1}$ est donnée \inserted{``en pointillés''} \uline{dans Cat hot}, de façon à rendre le carré commutatif,

% archive sheet number: 48
\scanpage{50}
\mspage{48}

alors on peut la réaliser par une flèche \uline{dans Cat} qui rende le diagramme comm. pour de bon (\unsure{dans} Cat). On est ramené au $=$ problème dans une \deleted{\unsure{situa.}}\unsure{situation}

\diagram{\begin{tikzcd}[ampersand replacement=\&]
X \arrow[r,hook,"i"'] \& Y \arrow[l,bend right=42,dashed,"p"']
\end{tikzcd}},

où $X$ est une catégorie de Thomason, $i$ une cofibration de Thomason, $p$ une rétraction de $Y$ dans $X$ \uline{dans} \uline{Cat hot}, on veut la réaliser par une vraie rétraction.

\margin{? $\big|$}

En admettant que ce soit le cas, donc que \deleted{\unsure{tout}} hom \deleted{\unsure{des}} ind objets $\underline{\underline K},\underline{\underline {K'}}\deleted{\unsure{{}_{\infty}}}$ \uline{dans Cat hot} se représente \unsure{par} un hom \deleted{de} entre les ind-objets \uline{dans Cat}, donc aussi par un foncteur $\deleted{\unsure{C}}K_\infty\to K'_\infty$, reste la question si cet homomorphisme \inserted{correspondant dans Cat hot} est unique. \deleted{dans le}

Donc si on a deux \deleted{\unsure{hom}} foncteurs $f,g$ qui induisent les $=$ hom sur les foncteurs ind-représentés \unsure{p.} sur Cat hot, il faut \unsure{q|que} \unsure{soient} homotopes. \inserted{Donc} \deleted{\unsure{on a donc à cons.}}

% slanted
\margin{$K,K'$ catégories de Thomason, les $K_n,K'_n$ aussi}

\uline{Question} Deux foncteurs $f,g\colon K_\infty\unsure{\relbar\joinrel\relbar|\longrightarrow} K'_\infty$ \deleted{\unsure{qui}} tels, que pour $\forall n\in\mathbb N$, \unsure{$f_n,g_n$|$f,g$} qui \unsure{s'en} déduisent par restriction à $K_n$ sont homotopes \unsure{?|,} (et $=$ homotopes en tant que foncteurs de $K_n$ dans un $K'_N$ convenable) \deleted{\unsure{que}} sont-ils homotopes ?

% archive sheet number: 49
